\documentclass[12pt,a4paper]{article}
\usepackage[utf8]{inputenc}
\usepackage[T1]{fontenc}
\usepackage[polish]{babel}
\usepackage{geometry}
\usepackage{booktabs}
\usepackage{array}
\usepackage{longtable}
\usepackage{hyperref}
\usepackage{xcolor}
\usepackage{listings}
\usepackage{enumitem}
\usepackage{graphicx}
\usepackage{float}

\geometry{margin=2.5cm}

\hypersetup{hidelinks}

\definecolor{codebg}{rgb}{0.95,0.95,0.95}
\definecolor{codeframe}{rgb}{0.7,0.7,0.7}

\lstset{
  backgroundcolor=\color{codebg},
  frame=single,
  rulecolor=\color{codeframe},
  basicstyle=\ttfamily\small,
  breaklines=true,
  captionpos=b,
  inputencoding=utf8,
  extendedchars=true,
  literate={ą}{{\k{a}}}1 {ę}{{\k{e}}}1 {ó}{{\'o}}1 {ś}{{\'s}}1
           {ł}{{\l{}}}1 {ż}{{\.z}}1 {ź}{{\'z}}1 {ć}{{\'c}}1
           {ń}{{\'n}}1 {Ą}{{\k{A}}}1 {Ę}{{\k{E}}}1 {Ó}{{\'O}}1
           {Ś}{{\'S}}1 {Ł}{{\L{}}}1 {Ż}{{\.Z}}1 {Ź}{{\'Z}}1
           {Ć}{{\'C}}1 {Ń}{{\'N}}1,
}

\title{\textbf{Dokumentacja funkcjonalna}\\\large Aplikacja Java -- Wizualizacja grafów planarnych}
\author{Mateusz Kamieniecki \\ Jan Rękas}
\date{Maj 2026}

\begin{document}

\maketitle
\tableofcontents
\newpage

% ---------------------------------------------------------------
\section{Cel projektu}

Celem projektu jest stworzenie aplikacji w~języku Java z~graficznym interfejsem użytkownika opartym na bibliotece Swing, umożliwiającej wczytywanie, wizualizację oraz interaktywną edycję układu wierzchołków grafów planarnych. Aplikacja wywołuje zewnętrzny binarny silnik obliczeniowy w~języku C (opracowany przez zespół Poliny Vyshni i~Andreia Udavichenki) w~celu wyznaczenia współrzędnych wierzchołków przy użyciu wybranego algorytmu, a~następnie prezentuje wyniki graficznie. Możliwe jest również pominięcie etapu obliczeń i~bezpośrednie wczytanie wcześniej przygotowanego pliku pozycji. Interfejs użytkownika jest w~całości w~języku polskim.

% ---------------------------------------------------------------
\section{Opis funkcjonalności}

Aplikacja udostępnia następujące funkcje:

\begin{itemize}
    \item wczytywanie opisu grafu (lista krawędzi) z~pliku tekstowego za pomocą okna dialogowego,
    \item wywołanie zewnętrznego binarnego silnika obliczeniowego jako procesu potomnego z~odpowiednimi argumentami (algorytm, ścieżka wejściowa, wyjściowa, format zapisu),
    \item wczytanie gotowego pliku pozycji bez uruchamiania obliczeń -- w~formacie tekstowym (.txt) lub binarnym (bez rozszerzenia),
    \item automatyczne nazwanie pliku wyjściowego na podstawie nazwy pliku wejściowego (przyrostek \texttt{\_positions}, zapis w~podkatalogu \texttt{positions} obok pliku grafu),
    \item wyświetlenie przykładowego grafu demonstracyjnego od razu po uruchomieniu aplikacji (4 wierzchołki, 4 krawędzie),
    \item wyświetlenie grafu w~panelu wizualizacji po zakończeniu obliczeń lub wczytaniu pliku pozycji,
    \item interaktywna edycja pozycji wierzchołka metodą przeciągnij i~upuść (ang.~\textit{drag \& drop}) bezpośrednio w~panelu wizualizacji,
    \item panowanie widokiem (przesuwanie środkowym przyciskiem myszy),
    \item zmiana skali widoku przy użyciu kółka myszy,
    \item przyciski pomocnicze: Reset zoom, Wyśrodkuj widok, Dopasuj do widoku,
    \item przełączanie widoczności identyfikatorów wierzchołków i~wag krawędzi,
    \item zapis zmodyfikowanych pozycji wierzchołków do pliku tekstowego,
    \item czytelne komunikaty o~błędach wyświetlane w~oknach dialogowych interfejsu graficznego.
\end{itemize}

% ---------------------------------------------------------------
\section{Interfejs graficzny}

Okno główne aplikacji podzielone jest na trzy obszary: lewy panel boczny zajmujący około 25--30\% szerokości okna (270\,px), prawy obszar wizualizacji grafu oraz wąski pasek stanu na dole okna. Interfejs nie posiada poziomego paska menu -- nawigacja między ekranami odbywa się za pomocą przycisku powrotu wbudowanego w~panel interakcji. Wszystkie napisy, etykiety, przyciski i~komunikaty są wyłącznie w~języku polskim.

\subsection{Lewy panel boczny}

Panel boczny wyświetla dwa różne ekrany przełączane automatycznie po wczytaniu grafu, zaimplementowane jako \texttt{CardLayout} z~kartami \texttt{settings} i~\texttt{interaction}.

% - - - - - - - - - - - - - - - - - - - - - -
\subsubsection{Karta: Ustawienia grafu (\texttt{settings})}

Karta aktywna domyślnie przy starcie aplikacji. Służy do skonfigurowania danych wejściowych, wyboru sposobu wyznaczania pozycji i~uruchomienia procesu. Elementy ułożone od góry do dołu:

\begin{enumerate}
    \item \textbf{Plik grafu (lista krawędzi)} -- pole tekstowe ze ścieżką do pliku oraz przycisk \ldots{} otwierający systemowe okno dialogowe wyboru pliku. Wartość domyślna: puste (użytkownik musi wskazać plik ręcznie). Akceptowany format: plik tekstowy z~listą krawędzi.

    \item \textbf{Tryb} -- cztery przyciski opcji (radio button) podzielone separatorem na dwie sekcje:
    
    \textit{Sekcja 1 -- Obliczanie (wywołuje silnik C):}
    \begin{itemize}
        \item \textbf{Fruchterman-Reingold} -- algorytm siłowy oparty na symulacji fizyki (opcja domyślna),
        \item \textbf{Triangulacja} -- algorytm geometryczny gwarantujący brak przecięć krawędzi.
    \end{itemize}
    Po wybraniu jednej z~tych opcji sekcja \textit{Format zapisu} staje się widoczna, a~sekcja \textit{Plik pozycji} jest ukryta.

    \textit{Sekcja 2 -- Wczytaj gotowy plik pozycji (pomija obliczenia):}
    \begin{itemize}
        \item \textbf{Wczytaj gotowy plik .txt} -- wczytuje plik tekstowy z~współrzędnymi,
        \item \textbf{Wczytaj gotowy plik binarny} -- wczytuje plik binarny.
    \end{itemize}
    Po wybraniu jednej z~opcji wczytywania sekcja \textit{Plik pozycji} staje się widoczna, a~sekcja \textit{Format zapisu} jest ukryta. Kliknięcie \textit{Wykonaj} wczytuje wskazany plik i~wyświetla graf bez wywoływania silnika C.

    \item \textbf{Plik pozycji} -- widoczny tylko w~trybie wczytywania. Pole tekstowe ze ścieżką do pliku pozycji oraz przycisk \ldots. Akceptowane formaty: \texttt{.txt} lub plik binarny (bez rozszerzenia).

    \item \textbf{Format zapisu} -- widoczny tylko w~trybie obliczania. Etykieta informuje, że plik zostanie zapisany w~podkatalogu \texttt{positions} obok pliku grafu. Dwa przyciski opcji (radio button):
    \begin{itemize}
        \item \texttt{.txt} -- zapis w~formacie tekstowym,
        \item \texttt{binarny (bez rozszerzenia)} -- zapis w~formacie binarnym (opcja domyślna).
    \end{itemize}
    Możliwy jest wybór tylko jednego formatu jednocześnie. Nazwa pliku wyjściowego tworzona jest automatycznie jako nazwa pliku wejściowego z~przyrostkiem \texttt{\_positions} (np.~\texttt{graf\_positions} lub \texttt{graf\_positions.txt}).

    \item \textbf{Przycisk \textit{Wykonaj}} -- działanie zależy od wybranego trybu:
    \begin{itemize}
        \item tryb obliczania: aplikacja wywołuje silnik C z~odpowiednimi argumentami; po zakończeniu wyświetla graf, aktualizuje pasek stanu i~przełącza panel na kartę interakcji,
        \item tryb wczytywania: aplikacja odczytuje wskazany plik i~wyświetla graf bez wywoływania silnika C; następnie przełącza na kartę interakcji.
    \end{itemize}
    Jeśli wymagane pole ścieżki jest puste lub plik nie istnieje, wyświetlany jest komunikat o~błędzie i~żadna akcja nie jest podejmowana.
\end{enumerate}

% - - - - - - - - - - - - - - - - - - - - - -
\subsubsection{Karta: Interakcja (\texttt{interaction})}

Karta staje się aktywna automatycznie po pomyślnym wczytaniu lub obliczeniu pozycji grafu. Elementy ułożone od góry do dołu:

\textit{Sekcja: Wyświetlanie}
\begin{itemize}
    \item \textbf{Identyfikatory wierzchołków} -- pole wyboru (check box); domyślnie włączone. Gdy aktywne, przy każdym wierzchołku wyświetlana jest jego etykieta z~identyfikatorem.
    \item \textbf{Wagi krawędzi} -- pole wyboru; domyślnie włączone. Gdy aktywne, przy każdej krawędzi wyświetlana jest jej nazwa i~waga.
\end{itemize}

\textit{Sekcja: Widok}
\begin{itemize}
    \item \textbf{Resetuj zoom} -- przywraca skalę 100\% (scale = 1.0).
    \item \textbf{Wyśrodkuj widok} -- resetuje przesunięcie widoku (panX = 0, panY = 0) bez zmiany skali.
    \item \textbf{Dopasuj do widoku} -- automatycznie dobiera skalę i~przesunięcie tak, by cały graf był widoczny w~panelu z~marginesem 80\,px.
\end{itemize}

\textit{Sekcja: Eksport}
\begin{itemize}
    \item \textbf{Zapisz zmiany pozycji} -- zapisuje bieżące współrzędne wszystkich wierzchołków do pliku tekstowego. Nazwa pliku: \texttt{<nazwa\_pliku\_grafu>\_newPositions.txt} w~tym samym katalogu co plik grafu.
\end{itemize}

\noindent Przycisk powrotu: \textit{$\leftarrow$ Wróć do ustawień} -- przełącza panel z~powrotem na kartę \textit{Ustawienia grafu}.

% ---------------------------------------------------------------
\subsection{Panel wizualizacji (prawy)}

Prawy obszar okna wypełnia niestandardowy komponent Swing (\texttt{GraphPanel}) z~własną implementacją metody \texttt{paintComponent()}. Po uruchomieniu aplikacji wyświetlany jest demonstracyjny graf składający się z~4 wierzchołków (id: 1--4) i~4 krawędzi (AB, BC, CD, DA), rozmieszczonych w~kwadracie 100$\times$100\,px. Graf użytkownika pojawia się po kliknięciu \textit{Wykonaj}.

Komponent rysuje:
\begin{itemize}
    \item krawędzie grafu jako odcinki proste łączące odpowiednie wierzchołki (kolor zielonkawy),
    \item opcjonalne etykiety nazwy i~wagi krawędzi w~punkcie środkowym krawędzi,
    \item wierzchołki jako wypełnione okręgi o~promieniu 12\,px (kolor niebieski),
    \item wierzchołek aktualnie przeciągany wyświetlany jest na czerwono,
    \item opcjonalne etykiety identyfikatorów wierzchołków (wyświetlane 15\,px nad środkiem okręgu).
\end{itemize}

Obsługiwane zdarzenia myszy:
\begin{itemize}
    \item Kółko myszy -- powiększanie/pomniejszanie (współczynnik 1.1 na każdy skok kółka).
    \item Środkowy przycisk myszy (przytrzymanie i~przeciągnięcie) -- panowanie widokiem.
    \item Lewy przycisk myszy (kliknięcie na wierzchołku i~przeciągnięcie) -- przesuwanie wierzchołka (drag \& drop).
\end{itemize}

% ---------------------------------------------------------------
\subsection{Pasek stanu}

Jednowierszowy pasek na dole okna (wysokość 28\,px) wyświetla bieżący stan programu, np.:
\begin{itemize}
    \item \texttt{Gotowy} -- stan domyślny przy starcie,
    \item \texttt{Obliczanie (fruchterman)\ldots} -- podczas pracy silnika C,
    \item \texttt{Wczytano pozycje z .txt} / \texttt{Wczytano pozycje binarne} -- po pomyślnym wczytaniu,
    \item \texttt{Gotowy --- N wierzchołków} -- po zakończeniu obliczeń,
    \item \texttt{Zoom zresetowany} / \texttt{Widok wyśrodkowany} / \texttt{Graf dopasowany do okna} -- po akcjach widoku,
    \item \texttt{Zapisano: <nazwa\_pliku>} -- po zapisaniu pozycji.
\end{itemize}

% ---------------------------------------------------------------
\section{Formaty danych}

Aplikacja korzysta z~formatów danych zdefiniowanych przez silnik obliczeniowy. Poniżej opisano każdy z~nich.

\subsection{Plik wejściowy -- lista krawędzi}

Plik tekstowy z~listą krawędzi grafu. Każda linia reprezentuje jedną krawędź:

\begin{lstlisting}
<nazwa_krawedzi> <wierzcholek_A> <wierzcholek_B> <waga_krawedzi>
\end{lstlisting}

\vspace{0.3cm}
\noindent \textbf{Przykład:}
\begin{lstlisting}
AB   1  2  1
BC   2  3  1
CD   3  4  1
DB   4  2  1.407
\end{lstlisting}

\noindent Gdzie:
\begin{itemize}
    \item \texttt{nazwa\_krawedzi} -- ciąg znaków bez spacji (np.~\texttt{E1}, \texttt{AB}),
    \item \texttt{wierzcholek\_A}, \texttt{wierzcholek\_B} -- liczby całkowite identyfikujące węzły,
    \item \texttt{waga\_krawedzi} -- liczba zmiennoprzecinkowa.
\end{itemize}

Linie puste są pomijane. Linia o~mniej niż 4 polach lub błędnym typie danych powoduje wyjątek \texttt{IOException} z~numerem linii.

% - - - - - - - - - - - - - - - - - - - - - -
\subsection{Plik wyjściowy -- format tekstowy (.txt)}

Plik zawiera listę węzłów z~obliczonymi współrzędnymi. Brak nagłówka. Każda linia to jeden wierzchołek:

\begin{lstlisting}
<vertex_id> <x> <y>
\end{lstlisting}

Współrzędne są zapisywane w~formacie \texttt{\%.2f} (dwa miejsca po przecinku). Przy wczytywaniu akceptowane są zarówno przecinek, jak i~kropka jako separator dziesiętny.

\vspace{0.3cm}
\noindent \textbf{Przykład:}
\begin{lstlisting}
1 0.00 0.00
2 1.00 0.00
3 1.00 1.00
4 0.00 1.00
\end{lstlisting}

% - - - - - - - - - - - - - - - - - - - - - -
\subsection{Plik wyjściowy -- format binarny (bez rozszerzenia)}

Plik binarny bez nagłówka, zapisywany w~kolejności little-endian. Każdy rekord ma stałą długość 20 bajtów:

\begin{center}
\renewcommand{\arraystretch}{1.3}
\begin{tabular}{>{\raggedright\arraybackslash}p{4.5cm} >{\raggedright\arraybackslash}p{7cm}}
\toprule
\textbf{Pole} & \textbf{Format} \\
\midrule
ID wierzchołka          & 4 bajty -- 32-bitowa liczba całkowita (little-endian) \\
Współrzędna $x$         & 8 bajtów -- IEEE 754 double (little-endian) \\
Współrzędna $y$         & 8 bajtów -- IEEE 754 double (little-endian) \\
\bottomrule
\end{tabular}
\end{center}

Rekordy następują bezpośrednio po sobie do końca pliku. Aplikacja Java odczytuje plik porcjami po 20 bajtów (\texttt{ByteBuffer}, porządek \texttt{LITTLE\_ENDIAN}).


% ---------------------------------------------------------------
\section{Dostępne algorytmy silnika obliczeniowego}

Silnik obliczeniowy udostępnia dwa algorytmy wybierane flagą \texttt{-a} przy wywołaniu:

\begin{center}
\renewcommand{\arraystretch}{1.5}
\begin{tabular}{>{\raggedright\arraybackslash}p{3.5cm} >{\raggedright\arraybackslash}p{10cm}}
\toprule
\textbf{Algorytm} & \textbf{Opis} \\
\midrule
\texttt{fruchterman} & Algorytm siłowy (Fruchterman-Reingold). Symulacja fizyczna: węzły odpychają się jak cząstki naładowane, krawędzie przyciągają je jak sprężyny. Iteracyjne chłodzenie zapewnia równomierny i~czytelny układ. Algorytm domyślny. \\
\addlinespace
\texttt{triangulation} & Algorytm geometryczny. Zapobiega krzyżowaniu się krawędzi przez precyzyjne rozmieszczanie wierzchołków z~użyciem reguł triangulacji. Graf musi być planarny; w~przeciwnym razie silnik zgłasza błąd. \\
\bottomrule
\end{tabular}
\end{center}

\vspace{0.3cm}
Przykładowe wywołanie silnika (generowane przez klasę \texttt{GraphSettingsCard}):
\begin{lstlisting}
engine.exe input.txt -a fruchterman -o positions/graf_positions -b
engine.exe input.txt -a triangulation -o positions/graf_positions.txt -t
\end{lstlisting}

Silnik wyszukiwany jest kolejno w~lokalizacjach: \texttt{bin/engine/engine.exe}, \texttt{app/bin/engine/engine.exe}, \texttt{../bin/engine/engine.exe}.

% ---------------------------------------------------------------
\section{Przykładowe scenariusze użycia}

\subsection{Obliczenie układu i wyświetlenie grafu}
\begin{enumerate}
    \item Uruchom aplikację. Panel wizualizacji wyświetla demonstracyjny graf (4 wierzchołki, 4 krawędzie).
    \item W~polu \textit{Plik grafu} kliknij \ldots{} i~wskaż plik z~listą krawędzi.
    \item Wybierz algorytm: \textit{Fruchterman-Reingold} lub \textit{Triangulacja}.
    \item Opcjonalnie zmień format zapisu (.txt lub binarny).
    \item Kliknij \textbf{Wykonaj}. Aplikacja wywołuje silnik C. Po zakończeniu obliczeń panel wizualizacji wyświetla graf, a~panel boczny przełącza się automatycznie na kartę interakcji.
\end{enumerate}

\subsection{Wczytanie gotowego pliku pozycji}
\begin{enumerate}
    \item W~karcie \textit{Ustawienia grafu} wybierz opcję \textit{Wczytaj gotowy plik .txt} lub \textit{Wczytaj gotowy plik binarny}.
    \item W~polu ścieżki kliknij \ldots{} i~wskaż plik pozycji.
    \item Wskaż również plik grafu (lista krawędzi) -- jest wymagany zawsze.
    \item Kliknij \textbf{Wykonaj}. Aplikacja odczytuje oba pliki i~wyświetla graf. Silnik C nie jest wywoływany.
\end{enumerate}

\subsection{Edycja pozycji wierzchołka (drag \& drop)}
\begin{enumerate}
    \item Po wyświetleniu grafu kliknij i~przytrzymaj lewy przycisk myszy na wybranym wierzchołku w~panelu wizualizacji.
    \item Przeciągnij wierzchołek w~nowe miejsce -- przesuwa się na bieżąco wraz z~kursorem; wierzchołek zmienia kolor na czerwony.
    \item Zwolnij przycisk myszy. Nowe współrzędne są zapisane w~modelu grafu.
\end{enumerate}

\subsection{Zapis zmodyfikowanych pozycji}
\begin{enumerate}
    \item Upewnij się, że graf jest wczytany (karta interakcji jest aktywna).
    \item Kliknij \textbf{Zapisz zmiany pozycji}.
    \item Aplikacja zapisuje współrzędne do pliku \texttt{<nazwa\_grafu>\_newPositions.txt} w~tym samym katalogu co plik grafu i~wyświetla komunikat potwierdzający.
\end{enumerate}

\subsection{Dopasowanie widoku}
\begin{enumerate}
    \item Kliknij \textbf{Dopasuj do widoku} -- aplikacja automatycznie dobiera skalę i~pozycję kamery tak, by cały graf był widoczny z~marginesem 80\,px.
    \item Użyj kółka myszy do ręcznej regulacji powiększenia.
    \item Przytrzymaj środkowy przycisk myszy i~przeciągnij, by przesunąć widok.
\end{enumerate}

% ---------------------------------------------------------------
\section{Sytuacje wyjątkowe}

Wszystkie komunikaty o~błędach wyświetlane są w~oknie dialogowym (\texttt{JOptionPane}). Aplikacja pozostaje uruchomiona we wszystkich poniższych przypadkach -- wczytany graf i~obliczone współrzędne nie są tracone.

\vspace{0.3cm}
\renewcommand{\arraystretch}{1.4}
\begin{longtable}{>{\raggedright\arraybackslash}p{6cm} >{\raggedright\arraybackslash}p{7.5cm}}
\toprule
\textbf{Sytuacja} & \textbf{Komunikat / zachowanie} \\
\midrule
\endfirsthead
\toprule
\textbf{Sytuacja} & \textbf{Komunikat / zachowanie} \\
\midrule
\endhead
\bottomrule
\endfoot
Nie podano pliku grafu & \texttt{Błąd: Nie podano pliku grafu.} \\
\addlinespace
Plik grafu nie istnieje na dysku & \texttt{Błąd: Plik grafu nie istnieje:}\\ \texttt{<\textit{ścieżka}>} \\
\addlinespace
Błąd parsowania pliku grafu & \texttt{Błąd wczytywania grafu:}\\ \texttt{<\textit{opis}>} \\
\addlinespace
Nie podano pliku pozycji (tryb wczytywania) & \texttt{Błąd: Nie podano pliku pozycji.} \\
\addlinespace
Błąd wczytywania pliku pozycji .txt & \texttt{Błąd wczytywania pozycji:}\\ \texttt{<\textit{opis}>} \\
\addlinespace
Błąd wczytywania pliku binarnego & \texttt{Błąd wczytywania pozycji:}\\ \texttt{<\textit{opis}>} \\
\addlinespace
Nie znaleziono silnika C (engine.exe) & \texttt{Nie znaleziono silnika obliczeniowego.}\\(z~listą oczekiwanych lokalizacji i~bieżącym katalogiem) \\
\addlinespace
Silnik C zakończył z~kodem błędu & \texttt{Silnik zakończył z~kodem błędu <nr>:}\\ \texttt{<\textit{wyjście silnika}>} \\
\addlinespace
Brak pliku wyjściowego po działaniu silnika & \texttt{Błąd operacji na pliku:}\\ \texttt{Brak pliku TXT/binarnego: <\textit{ścieżka}>} \\
\addlinespace
Błąd zapisu pozycji do pliku .txt & \texttt{Błąd zapisu:}\\ \texttt{<\textit{opis}>} \\
\addlinespace
Brak wczytanego grafu przy zapisie & \texttt{Brak wczytanego grafu.} \\
\addlinespace
Plik binarny pusty lub niekompletny & \texttt{Plik binarny jest pusty lub ma}\\ \texttt{nieprawidłowy format} \\
\addlinespace
Plik wejściowy z~wierzchołkami pusty & \texttt{Błąd: wczytany plik wierzchołków}\\ \texttt{jest pusty} \\
\addlinespace
Wczytany graf nie zawiera krawędzi & \texttt{Błąd: wczytany graf jest pusty} \\
\bottomrule
\end{longtable}

% ---------------------------------------------------------------

\end{document}